\clearpage\fancyhead[L]{\small PALLAS.}\hypertarget{native-359-lower}{}\label{native-359-lower}
\begin{center}\textbf{PALLAS.}\par\vspace{4pt}\hrule\vspace{4pt}\fontsize{10}{12}\selectfont Monatliche Correspondenz. Vol. XXVI, pp. 199\;..\;203. August 1812.\par\vspace{4pt}\hrule\vspace{4pt}\end{center}

\begin{flushright}Göttingen, 4 August 1812.\end{flushright}
\par\indent Some time ago I told you about the first sighting of \emph{Pallas} this year. I have made only a few observations of my own since then. Around the time of opposition the planet stayed too high above our observing site all night for us to make micrometric comparisons with our best instrument. Since it was only of the tenth magnitude, it was not visible at all with the mural quadrant. I therefore made do with observing it once around that time with a weaker instrument fitted with a ring micrometer. The resulting declination was rather unreliable.

\Needspace{11\baselineskip}\par\indent Here are all my observations:

\begin{center}
\fontsize{10}{12}\selectfont\setlength{\tabcolsep}{3pt}\renewcommand{\arraystretch}{1.15}\begin{tabular}{c|c|c|c}
1812 & Mean time & Right ascension & Northern decl. \\
\hline
April 9 & $11^{\mathrm u}\;32^{\mathrm m}\;27^{\mathrm s}$ & $269^\circ\;5'\;58''4$ & $16^\circ\;38'\;0''\!:\,:$ \\
May 2 & 10 \quad 36 \quad 6 & 268 \quad 42 \quad 5.3 & 21 \quad 4 \quad 31.7 \\
4 & 10 \quad 44 \quad 44 & 268 \quad 31 \quad 38.0 & 21 \quad 25 \quad 4.3 \\
June 7 & 11 \quad 13 \quad 51 & 262 \quad 50 \quad 51.1 & 25 \quad 2 \quad 56.6 \\
\end{tabular}
\end{center}

\Needspace{9\baselineskip}\par\indent Dr. \textsc{Olbers} had made the following two observations in Bremen:

\begin{center}
\fontsize{10}{12}\selectfont\setlength{\tabcolsep}{3pt}\renewcommand{\arraystretch}{1.15}\begin{tabular}{c|c|c|c}
1812 & Mean time & Right ascension & Northern decl. \\
\hline
April 3 & $11^{\mathrm u}\;26^{\mathrm m}\;44^{\mathrm s}$ & $268^\circ\;42'\;23''1$ & $15^\circ\;25'\;4''9$ \\
4 & 12 \quad 2 \quad 33 & 268 \quad 47 \quad 12.5 & 15 \quad 37 \quad 41.1 \\
\end{tabular}
\end{center}

\Needspace{12\baselineskip}\par\indent From Paris he was kind enough to send me the following four meridian observations by Dr. \textsc{Burkhardt}:

\begin{center}
\fontsize{10}{12}\selectfont\setlength{\tabcolsep}{3pt}\renewcommand{\arraystretch}{1.15}\begin{tabular}{c|c|c|c}
1812 & Mean time & Right ascension & Northern decl. \\
\hline
June 5 & $12^{\mathrm u}\;35^{\mathrm m}\;57^{\mathrm s}4$ & $263^\circ\;15'\;24''0$ & $24^\circ\;58'\;28''5$ \\
6 & 12 \quad 33 \quad 10.5 & 263 \quad 2 \quad 35. & 25 \quad 0 \quad 30. \\
8 & 12 \quad 21 \quad 37.0 & 262 \quad 37 \quad 6.5 & 25 \quad 3 \quad 57.7 \\
11 & 12 \quad 7 \quad 15.7 & 261 \quad 58 \quad 37.8 & 25 \quad 6 \quad 29.3 \\
\end{tabular}
\end{center}

\Needspace{20\baselineskip}\par\indent The second observation was reported to be less accurate. These are all the observations from this year that I have heard of. Mr. \textsc{Nicolai} compared them all with the elements and found the following agreement:
\SourceLeaf{360}\begin{center}
\emph{Difference.}
\end{center}

\begin{center}\fontsize{10}{12}\selectfont\setlength{\tabcolsep}{6pt}\renewcommand{\arraystretch}{1.15}\begin{tabular}{c|c|c|c}
1812 & R.A. & Decl. & Observer\\\hline
April 3 & $+30''3$ & $-\!-10''7$ & Olbers\\
4 & $+33.3$ & $-16.8$ & Olbers\\
9 & $+30.6$ &  & Gauss\\
May 2 & $+57.6$ & $-5.4$ & Gauss\\
4 & $+49.4$ & $+6.6$ & Gauss\\
\end{tabular}\end{center}
\begin{center}\fontsize{10}{12}\selectfont\setlength{\tabcolsep}{6pt}\renewcommand{\arraystretch}{1.15}\begin{tabular}{c|c|c|c}
1812 & R.A. & Decl. & Observer\\\hline
June 5 & $+51''4$ & $+6''1$ & Burckhardt\\
6 & $+59.4$ & $+11.7$ & Burckhardt\\
7 & $+56.7$ & $-34.0$ & Gauss\\
8 & $+59.6$ & $+2.4$ & Burckhardt\\
11 & $+62.0$ & $+8.2$ & Burckhardt\\
\end{tabular}\end{center}

\par\indent I attribute it chiefly to the influence of Saturn and Mars that the elements differ by 1 minute at this year’s opposition. In deriving them, the perturbations from Jupiter were taken into account as precisely as possible, and the elements fitted all seven of the first oppositions so closely. Mr. \textsc{Nicolai} has already made a preliminary calculation of the perturbations due to Saturn under my supervision. A second, more precise calculation will soon be made either by this exceptionally skilled and practiced calculator or by me. I will discuss its outcome in more detail in due course.

\par\indent Mr. \textsc{Nicolai} determined the opposition from the June observations as follows.

\begin{center}\fontsize{10}{12}\selectfont\setlength{\tabcolsep}{0pt}\renewcommand{\arraystretch}{1.15}\begin{tabular}{l}
1812 June 10. $3^{\mathrm u}\;2^{\mathrm m}\;48^{\mathrm s}$\\Mean time, Göttingen meridian
\end{tabular}\end{center}
\begin{center}\fontsize{10}{12}\selectfont\setlength{\tabcolsep}{2pt}\renewcommand{\arraystretch}{1.15}\begin{tabular}{lcr}
True longitude & $.\;.\;.\;.\;.\;.\;.\;.\;.\;.\;.\;.\;.\;.\;.$ & $259^\circ\;27'\;51''7$\\
True geocentric latitude & $.\;.\;.\;.\;.\;.\;.\;.\;.\;.$ & $48\;18\;26.1\;\mathrm{N}.$\\
\end{tabular}\end{center}

\Needspace{18\baselineskip}\par\indent In the meantime I have tried to see how closely all eight oppositions observed so far can still be brought into agreement when only Jupiter’s perturbations are taken into account. The following corrections must be made to the elements for that purpose:

\begin{center}\fontsize{9}{11}\selectfont\setlength{\tabcolsep}{2pt}\renewcommand{\arraystretch}{1.15}\begin{tabular}{lcr}
Correction to the epoch 1810 & $.\;.\;.\;.\;.\;.\;.\;.\;.\;.\;.$ & $-13''31$\\
to the daily mean motion & $.\;.\;.\;.\;.\;.\;.\;.\;.\;.$ & $-0.01200$\\
to the longitude of perihelion & $.\;.\;.\;.\;.\;.\;.\;.\;.\;.\;.\;.$ & $-17''57$\\
to the angle of eccentricity & $.\;.\;.\;.\;.\;.\;.\;.\;.\;.\;.\;.\;.$ & $-4.09$\\
to the logarithm of the semimajor axis & $.\;.\;.\;.\;.\;.\;,$ & $+0.0000045$\\
to the longitude of the node & $.\;.\;.\;.\;.\;.\;.\;.\;.\;.\;.\;.\;.\;.$ & $+0''80$\\
to the inclination of the orbit & $.\;.\;.\;.\;.\;.\;.\;.\;.\;.\;.\;.\;.\;.$ & $-1''21$\\
\end{tabular}\end{center}

\Needspace{16\baselineskip}\noindent With these corrections all eight oppositions agree as follows:

\begin{center}
\emph{Difference}
\end{center}

\begin{center}
\fontsize{10}{12}\selectfont\setlength{\tabcolsep}{6pt}\renewcommand{\arraystretch}{1.15}\begin{tabular}{c|c|c}
Opposition & Mean longitude & Helioc. latitude \\
\hline
1803 & $+5''8$ & $-2''8$ \\
1804 & $-12.9$ & $-9.7$ \\
1805 & $-15.2$ & $-0.7$ \\
1807 & $+8.3$ & $-5.8$ \\
1808 & $+16.0$ & $+11.7$ \\
1809 & $+13.0$ & $-1.5$ \\
1811 & $+11.7$ & $-9.5$ \\
1812 & $-25.9$ & $+6.1$ \\
\end{tabular}
\end{center}

\par\indent I will add only that where the numbers given here differ somewhat from those already printed elsewhere, this is the result of a later correction to the calculation.

\Needspace{13\baselineskip}\par\indent Mr. \textsc{Nicolai} has calculated next year’s opposition \emph{in advance} as follows:
\SourceLeaf{361}\begin{center}\fontsize{10}{12}\selectfont\setlength{\tabcolsep}{0pt}\renewcommand{\arraystretch}{1.15}\begin{tabular}{l}
Time 1813 Aug. 18\quad $8^{\mathrm u}\;50^{\mathrm m}\;5^{\mathrm s}$\\Mean time in Göttingen
\end{tabular}\end{center}
\begin{center}\fontsize{10}{12}\selectfont\setlength{\tabcolsep}{2pt}\renewcommand{\arraystretch}{1.15}\begin{tabular}{lcr}
True longitude & $.\;.\;.\;.\;.\;.\;.\;.\;.$ & $325^\circ\;24'\;17''9$\\
Geocentr. latitude & $.\;.\;.\;,\;.\;.\;.\;.$ & $24\;37\;23.8\;\mathrm{N}.$\\
\end{tabular}\end{center}

\par\indent I have also observed \emph{Juno} several times. It was extremely faint this year. In my next letter I will give a detailed report on the results inferred from those observations, together with an ephemeris for next year calculated by Mr. \textsc{Wachter}.
